\thispagestyle{plain}
\chapter*{Résumé}

Dans le cadre de l'amélioration continue des processus de conception assistée par ordinateur au sein de Segula Technologies Maroc, ce projet de fin d'études porte sur l'optimisation automatique de la géométrie des pièces mécaniques en utilisant des réseaux de neurones convolutifs (CNN) couplés à l'outil de contrôle qualité Q-Checker. L'objectif principal est de développer des macros CATIA automatisées capables de détecter les non-conformités géométriques et de proposer des corrections en s'appuyant sur des modèles d'intelligence artificielle.

La démarche adoptée suit la méthodologie DMAIC (Définir, Mesurer, Analyser, Innover, Contrôler), permettant une approche structurée et rigoureuse du problème. Les résultats obtenus montrent une réduction significative du temps de vérification manuelle et une amélioration de la conformité des pièces aux standards métier de l'industrie automobile.

Ce travail contribue à la digitalisation et à l'automatisation des processus de contrôle qualité dans le domaine de la conception mécanique, ouvrant ainsi la voie à une intégration plus large de l'intelligence artificielle dans les flux de travail CAO.

\keywordspt{Optimisation géométrique, Réseaux de neurones convolutifs, Q-Checker, CATIA, DMAIC, CAO, Contrôle qualité.}

\MediaOptionLogicBlank

\pdfbookmark[1]{Abstract}{abstract}
\chapter*{Abstract}

As part of the continuous improvement of computer-aided design processes at Segula Technologies Morocco, this final year project focuses on the automatic optimization of mechanical part geometry using Convolutional Neural Networks (CNN) combined with the Q-Checker quality control tool. The main objective is to develop automated CATIA macros capable of detecting geometric non-conformities and suggesting corrections based on artificial intelligence models.

The approach follows the DMAIC methodology (Define, Measure, Analyze, Improve, Control), enabling a structured and rigorous problem-solving framework. The results demonstrate a significant reduction in manual verification time and an improvement in part compliance with automotive industry standards.

This work contributes to the digitalization and automation of quality control processes in the field of mechanical design, paving the way for a broader integration of artificial intelligence into CAD workflows.

\keywordsen{Geometry Optimization, Convolutional Neural Networks, Q-Checker, CATIA, DMAIC, CAD, Quality Control.}

\MediaOptionLogicBlank